\title{The Era of 1-bit LLMs: All Large Language Models are in 1.58 Bits}

\begin{abstract}
Advancements like BitNet~\cite{bitnet} have initiated a transformative phase in the realm of 1-bit Large Language Models (LLMs). In this study, we present \textbf{\text{BitNet b1.58}}, a novel 1-bit LLM configuration, wherein every model parameter (or weight) is restricted to ternary values from the set \{-1, 0, 1\}. Despite the drastic reduction in representation precision, this approach achieves parity with traditional full-precision (FP16 or BF16) Transformer LLMs of identical model scale and training regime in terms of perplexity as well as performance on downstream tasks. At the same time, it brings significant improvements in efficiency—reducing latency, memory consumption, throughput bottlenecks, and energy usage. Beyond empirical gains, this 1.58-bit LLM establishes a new scaling law, offering guidance for the development of future generations of LLMs that balance strong performance with cost efficiency. In addition, it paves the way for new computational paradigms and motivates the design of specialized hardware tailored for 1-bit LLM operations.
\end{abstract}

\section{The Era of 1-bit LLMs}

Over recent years, artificial intelligence has witnessed an exponential expansion in both the scale and capabilities of Large Language Models (LLMs). Although these systems excel across diverse natural language processing applications, the surge in their size intensifies deployment challenges and raises critical environmental and economic concerns, chiefly due to their considerable energy demands.

To tackle these issues, one prominent strategy is the use of post-training quantization for reducing bitwidth in models destined for inference~\cite{smoothquant, gptq, quip, quip_sharp}. By lowering the precision of parameters and activations, this methodology offers substantial savings in memory footprint and computational costs associated with LLMs. There has been ongoing progress, pushing quantization from the standard 16 bits down to formats as low as 4 bits~\cite{gptq, awq}. Nevertheless, despite its widespread deployment in practical LLM scenarios, post-training quantization remains a sub-optimal solution.

Recent advancements in 1-bit neural network architectures, exemplified by BitNet~\cite{bitnet}, offer a compelling pathway to decrease the operational costs of LLMs without sacrificing their accuracy. Standard LLMs utilize 16-bit floating-point formats (FP16 or BF16), with matrix multiplications constituting the primary workload. As a result, most computational expense is attributed to floating-point additions and multiplications. In comparison, BitNet's matrix operations rely solely on integer additions, leading to substantial reductions in energy consumption for LLMs. Since power constraints often limit computational performance in modern chips, such energy efficiencies also facilitate accelerated computation.

Beyond computation, transferring model parameters from DRAM to on-chip memory (like SRAM) can be resource-intensive during inference. While some strategies have sought to expand SRAM capacity to increase throughput, this comes at a much higher cost relative to DRAM. In contrast, 1-bit LLMs demand significantly less memory in terms of both storage capacity and bandwidth, compared to their full-precision counterparts. This markedly diminishes the time and expense involved in loading parameters from DRAM, thereby enabling quicker and more cost-effective inference.

In this paper, we propose an enhanced 1-bit LLM variant, denoted as \textbf{\text{BitNet b1.58}}, in which each model parameter can assume one of three discrete values: \{-1, 0, 1\}. This model extends the original 1-bit BitNet by introducing 0 as an allowable weight, equating to 1.58 bits per parameter in binary terms. \text{BitNet b1.58} preserves all the computational advantages of BitNet, including a paradigm shift towards matrix multiplication that nearly eliminates multiplications and can be aggressively optimized. Energy requirements remain as low as those for the original BitNet, while memory footprint, throughput, and latency greatly surpass those of FP16 LLMs. Moreover, \text{BitNet b1.58} provides two further benefits. Firstly, it exhibits enhanced modeling capacity by explicitly enabling feature selection—an improvement made possible through the 0-valued weights, leading to marked performance gains for 1-bit LLMs. Secondly, empirical evidence demonstrates that \text{BitNet b1.58} can rival full-precision FP16 models in terms of both perplexity and downstream task outcomes, beginning from a 3B parameter scale when all other factors (model size, training data, etc.) remain constant.

\section{BitNet b1.58}

\text{BitNet b1.58} is an extension of the BitNet framework, wherein the standard \emph{nn.Linear} layers are substituted with \emph{BitLinear} layers in the Transformer model. This variant is trained from the beginning with weights quantized to 1.58 bits and activations quantized to 8 bits. Relative to its predecessor, BitNet, several adjustments have been incorporated, as outlined below.

\paragraph{Quantization Function.} In order to restrict the weight values to the discrete set $\{-1, 0, +1\}$, we utilize an \emph{absmean} quantization scheme. This method normalizes the weight matrix by its mean absolute value prior to quantization, then rounds each resulting element to the nearest value in $\{-1, 0, +1\}$ as follows:

\begin{equation}
    \widetilde{W} = \mathrm{RoundClip}(\frac{W}{\gamma + \epsilon}, -1, 1),
\end{equation}

\begin{equation}
    \mathrm{RoundClip}(x, a, b) = \max(a, \min(b, \mathrm{round}(x))),
\end{equation}

\begin{equation}
    \gamma = \frac{1}{nm}\sum_{ij} |W_{ij}|.
\end{equation}

For activation quantization, we adhere to the original BitNet methodology with one exception: activations are not scaled into the interval $[0, Q_b]$ prior to the application of non-linearities. Instead, per-token scaling to $[-Q_b, Q_b]$ is applied, thus obviating the need for zero-point quantization. This strategy simplifies implementation and streamlines system-level optimizations, while causing only minor impact on model performance in our evaluation.

\paragraph{LLaMA-alike Components.}

LLaMA~\cite{llama, llama2} has become the standard architectural choice for open-source large language models (LLMs). To better align with this ecosystem, \text{BitNet b1.58} incorporates LLaMA-inspired architectural choices. Namely, the model adopts RMSNorm~\cite{rmsnorm}, SwiGLU~\cite{swiglu}, and rotary embedding~\cite{rope}, and omits all bias terms. These modifications facilitate seamless integration of \text{BitNet b1.58} into widely-used open-source frameworks, such as Huggingface, vLLM~\cite{vllm}, and llama.cpp\footnote{\url{https://github.com/ggerganov/llama.cpp}}, with minimal required adaptation.

\section{Results}

We conducted a comparative analysis between \text{BitNet b1.58} and our reproduced FP16 \text{LLaMA LLM} across different model scales. To maintain experimental consistency, both models underwent pre-training on the RedPajama dataset~\cite{redpajama} for a total of 100 billion tokens. We assessed their zero-shot capabilities over several language benchmarks, including ARC-Easy~\cite{arc}, ARC-Challenge~\cite{arc}, Hellaswag~\cite{hellaswag}, Winogrande~\cite{winoGrande}, PIQA~\cite{piqa}, OpenbookQA~\cite{openbookqa}, and BoolQ~\cite{boolq}. Additionally, we measured validation perplexity using the WikiText2~\cite{wikitext2} and C4~\cite{c4} datasets.

Furthermore, we evaluated and contrasted the runtime GPU memory consumption and inference latency of both \text{LLaMA LLM} and \text{BitNet b1.58}. For benchmarking, we employed the FasterTransformer\footnote{\url{https://github.com/NVIDIA/FasterTransformer}} framework, recognized for its GPU inference optimization for LLMs. For \text{BitNet b1.58}, we integrated the 2-bit kernel from Ladder~\cite{ladder}. The primary metric reported was time per output token, given its predominant impact on inference cost.

In Table~\ref{tab:ppl}, the results concerning perplexity and computational cost are presented for both \text{BitNet b1.58} and \text{LLaMA LLM}. Notably, from the 3B parameter scale onwards, \text{BitNet b1.58} achieves parity with full-precision \text{LLaMA LLM} in perplexity, while attaining a 2.71$\times$ speedup and a 3.55$\times$ reduction in GPU memory usage. Particularly, the 3.9B \text{BitNet b1.58} configuration is observed to be 2.4$\times$ faster and 3.32$\times$ more memory-efficient, all while significantly outperforming the 3B \text{LLaMA LLM} variant.

A comprehensive breakdown of zero-shot accuracy results is given in Table~\ref{tab:zero-shot}, which are derived using the evaluation protocol provided by \emph{lm-evaluation-harness}\footnote{\url{https://github.com/EleutherAI/lm-evaluation-harness}}. The findings indicate that as model scale grows, the performance differential between \text{BitNet b1.58} and \text{LLaMA LLM} diminishes. Crucially, \text{BitNet b1.58} matches the accuracy of the full-precision baseline beginning at the 3B model size. Consistent with the perplexity trends, the downstream results demonstrate that \text{BitNet b1.58} at 3.9B parameters surpasses \text{LLaMA LLM} 3B in both performance and efficiency, requiring less memory and exhibiting lower latency. These results position \text{BitNet b1.58} as a Pareto-optimal advancement over current leading LLMs.

\paragraph{Memory and Latency}

We extended our analysis to larger model configurations—specifically 7B, 13B, and 70B parameters—and assessed their associated costs. Figure~\ref{fig:latency-memory-energy} depicts the behavior of latency and memory usage, demonstrating that performance acceleration becomes more pronounced as models grow in scale. Notably, the \text{BitNet b1.58} 70B model achieves a 4.1-fold increase in speed relative to the \text{LLaMA LLM} reference. This improvement arises from the scaling of \emph{nn.Linear} computational overhead with model size. Memory utilization exhibits parallel dynamics: since embeddings are kept at full precision, their relative memory share diminishes as models become larger. All latency and memory figures are derived using a 2-bit kernel, suggesting further gains are possible through additional optimization.

\paragraph{Energy}

We provide an estimate of the energy requirements for arithmetic operations in both \text{BitNet b1.58} and \text{LLaMA LLM}. Our analysis emphasizes matrix multiplication, as it is the dominant contributor to the computational cost in large language models. As shown in Figure~\ref{fig:energy}, energy usage is mainly driven by INT8 addition in \text{BitNet b1.58}, while \text{LLaMA LLM} predominantly involves FP16 additions and multiplications. Based on the energy models from~\cite{energycost, pokebnn}, \text{BitNet b1.58} reduces matrix multiplication energy requirements by a factor of 71.4 when deployed on 7nm hardware. Additionally, we recorded end-to-end energy usage for processing 512 tokens. The findings reveal that energy efficiency advantages for \text{BitNet b1.58} over the FP16 \text{LLaMA LLM} baseline become more substantial as model size increases. This effect is explained by the growing share of \emph{nn.Linear} computations in larger architectures, while other model components contribute less to overall cost.

\paragraph{Throughput}

We evaluated and contrasted the throughput capabilities of \text{BitNet b1.58} and \text{LLaMA LLM} models, each with 70B parameters, deployed across two 80GB A100 GPUs. Pipeline parallelism~\cite{gpipe} was employed to enable execution of the \text{LLaMA LLM} 70B model on these GPUs. The batch size was incrementally increased until the memory limit of the GPUs was hit, keeping the sequence length fixed at 512. As summarized in Table~\ref{tab:throughput}, the 70B-parameter \text{BitNet b1.58} model accommodates batch sizes up to 11 times greater than those supported by \text{LLaMA LLM}, translating into a throughput advantage of 8.9 times.

\textbf{\text{BitNet b1.58} introduces a novel scaling paradigm relating model efficiency and inference resource cost.} According to our analysis, the following size equivalence between 1.58-bit and traditional 16-bit models can be drawn from Figures~\ref{fig:latency-memory-energy} and \ref{fig:energy}:

\begin{itemize}
    \item 13B BitNet b1.58 surpasses a 3B FP16 LLM in latency, memory footprint, and energy requirements.
    \item 30B BitNet b1.58 outperforms a 7B FP16 LLM across latency, memory utilization, and energy metrics.
    \item 70B BitNet b1.58 demonstrates greater efficiency in latency, memory, and energy than a 13B FP16 LLM.
\end{itemize}

\paragraph{Training with 2T Tokens}

The scale of training tokens is pivotal for LLM performance. To assess the token scalability of \text{BitNet b1.58}, we trained it with 2T tokens, adopting the same data preparation methodology used by StableLM-3B~\cite{StableLM-3B-4E1T}, which is the leading open-source 3B model. Both models were tested using a suite of benchmarks—Winogrande~\cite{winoGrande}, PIQA~\cite{piqa}, SciQ~\cite{sciq}, LAMBADA~\cite{lambada}, and ARC-easy~\cite{arc}—and zero-shot performance results are compiled in Table~\ref{tab:2t}. For tasks evaluated with both raw and normalized accuracy, the mean was reported. StableLM 3b's results at the 2T token scale were taken directly from its technical report. Our analysis reveals that \text{BitNet b1.58} attains higher performance on all tested tasks, underscoring the robust generalization abilities of 1.58-bit LLMs.

\section{Discussion and Future Work}

\textbf{1-bit Mixture-of-Experts (MoE) LLMs}

Mixture-of-Experts (MoE) architectures have established themselves as an efficient means for scaling LLMs while maintaining manageable computational demands. Although these models substantially lower the required floating-point operations, their practical utility is often constrained by high memory requirements and communication costs between chips. The adoption of 1.58-bit LLMs offers a promising solution to these bottlenecks. Reduced memory consumption enables the deployment of MoE models with fewer hardware resources and also diminishes network-related activation transfer overhead. In the ideal case where the entire model resides on a single chip, such overhead can be nearly eradicated.

\textbf{Native Support of Long Sequence in LLMs}

As the needs of contemporary LLMs increasingly focus on processing extended sequences, efficiently handling memory—especially for key-value (KV) cache storage—emerges as a primary obstacle for inference over long contexts. \text{BitNet b1.58} marks notable progress in supporting longer sequences natively, cutting activation precision from 16 bits down to 8 bits and effectively doubling the maximum attainable context length within the same resource constraints. There is potential for even more aggressive compression, losslessly reducing activations to 4 bits or less within 1.58-bit LLMs; we consider this an avenue for future investigation.

\textbf{LLMs on Edge and Mobile}

Deploying LLMs quantized to 1.58 bits can significantly boost the efficiency and performance of language models on mobile and edge platforms, which are frequently limited by restricted memory and computing resources. This level of quantization reduces both the energy and memory needs of LLMs, broadening the scope of applications that can run on such constrained devices and thereby unlocking new usage scenarios. Additionally, 1.58-bit LLMs offer improved compatibility with CPU-based systems, the prevalent processors in edge and mobile environments. As a result, models like \text{BitNet b1.58} can run efficiently in these contexts, further enhancing device functionality.

\textbf{New Hardware for 1-bit LLMs}

Innovations such as Groq\footnote{\url{https://groq.com/}} have demonstrated the promise of designing specialized hardware, including Language Processing Units (LPUs), for accelerating LLM workloads. Looking forward, we advocate for further research and development of hardware platforms and architectures tailored explicitly for 1-bit LLMs, taking advantage of the distinct computational properties enabled by frameworks such as BitNet~\cite{bitnet}.